\scanpage{20}
\mspage{19}

Or on a

\[
\deleted{\det}\Phi_{\mathscr L}(L\otimes L_\beta)=\Phi_{\mathscr L}(L)*\varepsilon_{-\beta}
\]

d'où

\[
\det\Phi_{\mathscr L}(L\unsure{\otimes|{}}L_\beta)\simeq\deleted{\unsure{\Phi}}\det\Phi_{\mathscr L}(L)*\varepsilon_{-\beta}\deleted{\unsure{=}}
\]

donc la classe de $\det(\Phi_{\mathscr L}(L\otimes L_\beta))$ dans $\underline{\mathrm{Pic}}_{B/S}(S)$ est égale à \deleted{$\det($} celle de $\det\Phi_{\mathscr L}(L)+\widetilde{\varphi(\delta)}(-\beta)$
% slanted

\margin{Vérif, avec Attention au \unsure{signe|signes}\ldots.}
(où $\varphi(\delta)$ \unsure{c'|est} l'image de $\det\Phi_{\mathscr L}(L)$ dans $\mathrm{NS}(B)$), d'où la conclusion annoncée.

Par suite, on peut interpréter $l_g$ comme un hom. $\underline{\mathrm{Pic}}^\delta_{X/S}\overset{B}{\times}(A,-\widetilde{\varphi(\delta)})$ dans $\underline{\mathrm{Pic}}^{\varphi(\delta)}_{B/S}$, ou encore comme une \deleted{\unsure{hom.}} section de $\underline{\mathrm{Pic}}^\delta_{X/S}\overset{B}{\times}(A,\widetilde{\varphi(\delta)})\overset{A}{\times}\underline{\mathrm{Pic}}^{\varphi(\delta)}_{B/S}$. On trouve ainsi \deleted{\unsure{$N$}}\inserted{un} hom (pour $g$ variable)

\[
X\mathbin{\unsure{-|\to}}\underline{\mathrm{Pic}}^{\varphi(\delta)}_{B/S}\overset{A}{\times}\underline{\mathrm{Pic}}^\delta_{X/S}\overset{B}{\times}(A,\unsure{\overline{\varphi(\delta)}|\widetilde{\varphi(\delta)}})
\]

qui n'est autre que \deleted{celui de} l'hom $l_\delta$ de d), comme on vérifie facilement.

\uline{Remarques}\qquad a) On n'a pas eu, dans cette interprétation, à faire intervenir un faisceau sur $\underline{\mathrm{Pic}}^\delta_{X/S}$, -- seult sur $B$.

Notons que si $\underline{\mathrm{Pic}}^0_{X/S}$ n'est pas un schéma abélien, on peut encore donner un sens aux constructions précédentes, en supposant \uline{donné} un schéma abélien
